\section{激活扰动几何}

本节把上一节的一阶更新几何转换为 activation perturbation 的语言。

考虑单层线性映射
\[
Z=WA,
\]
其中 \(A\) 是输入激活矩阵。更新 \(W^+=W-\eta D\) 后，直接 pre-activation 扰动为
\[
\Delta Z=-\eta DA.
\]

\subsection{Polar 更新的 operator-norm 控制}

若
\[
D=D_{\mathrm{polar}}=UV^\top,
\]
则
\[
\|D\|_{\op}=1.
\]
因此对任意单样本 activation vector \(a\)，有
\[
\|Da\|_2\le \|a\|_2.
\]
若步长为 \(\eta\)，则
\[
\|\Delta z(x)\|_2\le \eta\|h_{\ell-1}(x)\|_2.
\]
这说明 polar/SpecGrad 更新给出了 per-sample worst-case activation perturbation bound。

对整个 batch，有
\[
\|DA\|_F\le \|D\|_{\op}\|A\|_F=\|A\|_F.
\]
相反，Frobenius-normalized gradient direction
\[
D_F=G/\|G\|_F
\]
满足
\[
\|D_FA\|_F\le \|D_F\|_F\|A\|_{\op}=\|A\|_{\op}.
\]
由于
\[
\|A\|_F=\sqrt{\srank(A)}\|A\|_{\op},
\]
高稳定秩 activation 下，polar 更新可能导致更大的 batch-level activation movement。

因此，本文区分两种稳定性：
\[
\boxed{\text{operator/per-sample activation stability}}
\]
与
\[
\boxed{\text{batch-Frobenius activation stability}.}
\]
SpecGrad 更自然地控制前者，而非后者。

\subsection{层级谱更新优势条件}

设 loss 对 \(Z=WA\) 是 \(\beta\)-smooth。则
\[
\Loss(W-\eta D)
\le
\Loss(W)-\eta\langle G,D\rangle+\frac{\beta\eta^2}{2}\|DA\|_F^2.
\]

对 Frobenius-normalized direction \(D_F=G/\|G\|_F\)，有
\[
\langle G,D_F\rangle=\|G\|_F,
\qquad
\|D_FA\|_F\le \|A\|_{\op}.
\]
因此最优步长下的保证下降为
\[
\operatorname{Dec}_F\gtrsim
\frac{\|G\|_F^2}{2\beta\|A\|_{\op}^2}.
\]

对 polar direction \(D_S=UV^\top\)，有
\[
\langle G,D_S\rangle=\|G\|_*,
\qquad
\|D_SA\|_F\le \|A\|_F.
\]
因此
\[
\operatorname{Dec}_S\gtrsim
\frac{\|G\|_*^2}{2\beta\|A\|_F^2}.
\]

于是 spectral update 更优的充分条件为
\[
\frac{\|G\|_*^2}{\|G\|_F^2}
>
\frac{\|A\|_F^2}{\|A\|_{\op}^2},
\]
也就是
\[
\boxed{\nrank(G)>\srank(A).}
\]
